\section{Random sliding outer low-weight notes}
\label{sec:random-sliding-outer-low-weight-notes}

This note preserves fixed-tap low-weight facts for the random sliding dense outer. They are useful for finite scalar
diagnostics and for comparing sliding-band and block-code outer spectra, but they are not consumed by the current
compiled dense+dense asymptotic theorem.

\begin{lemma}[Fixed-tap outer removes weight one]
\label{lem:outer-fixedtap-no-weight-one}
For the fixed-tap systematic banded outer,
\[
\enum_{\Cone,1}=0
\]
deterministically.
\end{lemma}

\begin{proof}
A nonzero systematic outer codeword of weight one would have to come from a weight-one message \(x=e_i\) whose parity
part is zero. But the fixed diagonal tap gives \(p_i=1\), so the parity part has weight at least one.
\end{proof}

\begin{lemma}[Fixed-tap banded low-weight parity law]
\label{lem:outer-fixedtap-banded-lowweight-law}
Fix a message word \(x\) of weight \(w\). Let \(q(x)\) be the size of the union of the parity-band intervals touched by
the support of \(x\), and let \(r(x)\) be the number of band-clusters of that support, where adjacent support positions
are in the same cluster when their separation is at most \(M\). Then the fixed-tap parity weight
has the form
\[
r(x)+\Binomial(q(x)-r(x),\tfrac12).
\]
\end{lemma}

\begin{proof}
The first support position in each \(\sigma\)-cluster has a parity coordinate where only its fixed diagonal tap is
active, so that coordinate contributes a deterministic one. These give \(r(x)\) parity ones. Every other active parity
coordinate sees at least one random tap from the band, so its output is an independent fair bit. There are
\(q(x)-r(x)\) such coordinates.
\end{proof}
